\section{运动学C\quad\tiny 11.2-11.8}
\lsn{1}根据\(v - t\)图像得出匀变速直线运动的公式（至少7个）：

\begin{tikzpicture}[>=stealth,font=\footnotesize]
    \draw[|<->|,thin,gray] (3.25,2.5)--node[fill=white]{ \(\Delta v\)}(3.25,1);
    \filldraw[pattern=north east lines,draw=white,pattern color=red!30] (0,1)--(3,2.5)--(3,1)--(0,1);
    \filldraw[pattern=north west lines,draw=white,pattern color=red!30] (0,0) rectangle(3,1);
    \draw[->,blue] (0,0)--(5,0)node[below]{ \(t\)};
    \draw[->,blue] (0,0)node[below left]{\(O\)}--(0,4)node[right]{ \(v\)};
    \draw[magenta] (0,1)--(4,3)node[below right,black]{ \(a\)};
    \draw[dashed,cyan] (0,2.5)--(3,2.5)--(3,0) (0,1)--(3,1) (0,1.75)--(1.5,1.75)--(1.5,0);
    \node[blue] at(0,1)[left]{ \(v_0\)};
    \node[blue] at(0,2.5)[left]{ \(v_t\)};
    \node[blue] at(3,0)[below]{ \(t\)};
    \node[blue] at(0,1.75)[left]{ \(\overline{v}\)};
    \node[blue] at(1.5,0)[below]{\(\sfrac{t}{2}\)};
    \draw[|<->|,gray,thin] (0,-.8)--node[fill=white]{ \(\Delta t\)}(3,-.8);
  \end{tikzpicture}
\vspace*{2cm}



\lsn{2}表示出下列图像中运动的初速度\(v_0\)与加速度\(a\)，并给出过程

\begin{tikzpicture}[>=stealth,font=\footnotesize]
  \draw[->,blue] (0,0)node[below left]{\(O\)}--(4,0)node[below]{\(t\)};
  \draw[->,blue] (0,0)--(0,3)node[right]{\(\overline{v}\)};
  \draw[magenta] (0,.8)node[left,blue]{\(b\)}--(3,2.5)node[below,blue]{\(k\)};
  \begin{scope}[xshift=5cm]
    \draw[->,blue] (0,0)node[below left]{\(O\)}--(4,0)node[below]{\(\frac{1}{t}\)};
  \draw[->,blue] (0,0)--(0,3)node[right]{\(\frac{x}{t^2}\)};
  \draw[magenta] (0,.8)node[left,blue]{\(b\)}--(3,2.5)node[below,blue]{\(k\)};
  \end{scope}
  \begin{scope}[xshift=10cm]
    \draw[->,blue] (0,0)node[below left]{\(O\)}--(4,0)node[below]{\(x\)};
  \draw[->,blue] (0,0)--(0,3)node[right]{\(v^2\)};
  \draw[magenta] (0,.8)node[left,blue]{\(b\)}--(3,2.5)node[below,blue]{\(k\)};
  \end{scope}
\end{tikzpicture}

\vspace*{3cm}
\lsn{3}\(A\)、\(B\)两辆车（可视为质点）在一条直线上同向运动，\(A\)在\(B\)的正前方\(x_0\)处，它们的\(v - t\)图像如图所示，分析下列情况\(A\)、\(B\)是否相遇，相遇几次

\begin{tikzpicture}[>=stealth,font=\scriptsize]
  \draw[->,blue] (0,0)node[below left]{\(O\)}--(4,0)node[below]{\(t\)};
  \filldraw[pattern=north east lines,draw=white,pattern color=red!30] (0,.5)--(1.5,1.5)--(0,2);
  \draw[->,blue] (0,0)--(0,3)node[right]{\(v\)};
  \draw[magenta] (0,.5)--(3,2.5)node[right]{\(A\)};
  \draw[teal] (0,2)--(3,1)node[right]{\(B\)};
  \node at(.45,1.3){\(\Delta x\)};
\end{tikzpicture}
\begin{itemize}
  \item \(x_0 > \Delta x\)：
  \item \(x_0 = \Delta x\)：
  \item \(x_0 < \Delta x\)：
\end{itemize}

\lsn{4}某同学在实验中得到如下图所示的纸带，其中两点间有4个点未画出，打点计时器频率为\qty{50}{Hz}，计算并填空（保留两位有效数字）

\begin{tikzpicture}[>=stealth,font=\scriptsize]
	\coordinate (o) at(0,0);
	\draw ($(o)+(-.3,.5)$)--++(10.2,0) ($(o)+(-.3,-.5)$)--++(10.2,0)coordinate(a) ($(o)+(-.3,.5)$)--++(-.1,-.1)--++(.2,-.7)--++(-.1,-.2) (a)--++(.1,.1)--++(-.2,.7)--++(.1,.2);
	\node at($(a)+(-.7,-.2)$){单位\qty{}{cm}};
	
	\foreach \x/\y/\n in {0/.6/2.19,.6/1.6/2.76,1.6/3/3.35,3/4.8/3.93,4.8/7/4.58,7/9.6/5.20}
	  \draw[<->,thin] (\x,.1)--node[fill=white]{\tiny\(\n\)}(\y,.1);
  \foreach \x/\n in {0/0,.6/1,1.6/2,3/3,4.8/4,7/5,9.6/6}{
		\coordinate (\n) at(\x,0);
		\draw[gray] (\x,.05)--++(0,.15);
		\fill (\n)node[below]{\(\n\)} circle(.03);}
\end{tikzpicture}%%始.6,差.4


\vspace*{.5cm}
\(v_2 =\)\lunl{3cm}\qquad \(v_5 =\)\lunl{3cm}

\vspace*{.5cm}
\(a =\)\lunl{3cm}\qquad \(v_6 =\)\lunl{3cm}

\lsn{5}\icode{E-MK-SI-L2-260929-003}
汽车从甲城以速度大小\(v_1\)沿直线一直行驶到乙城，紧接着又从乙城以速度大小\(v_2\)沿直线返回，到达甲、乙两城中点的丙小镇，关于汽车在这一全过程的平均速度，下列说法中正确的是\paren[D]
\begin{choices}
	\item \(\dfrac{v_1 + v_2}{2}\)，方向为甲指向丙
	\item \(\dfrac{v_1 + v_2}{2}\)，方向为乙指向丙
	\item \(\dfrac{v_1v_2}{v_1 + 2v_2}\)，方向为乙指向丙
	\item \(\dfrac{v_1v_2}{v_1 + 2v_2}\)，方向为甲指向丙
\end{choices}
\begin{answer}
	D\qquad 提示：平均速度是总位移比总时间，注意汽车的运动方向变化
\end{answer}

\lsn{6}\icode{N16H-MK-SI-L3-261001-001}
\nceeID[3]{2016上海}{14}
物体做匀加速直线运动，相继经过两段距离为\qty{16}{m}的路程，第一段用时\qty{4}{s}，第二段用时\qty{2}{s}，则物体的加速度是\paren[B]
\begin{choices}
    \item \(\frac{2}{3}\) \qty{}{m/s^2}
    \item \(\frac{4}{3}\) \qty{}{m/s^2}
    \item \(\frac{8}{9}\) \qty{}{m/s^2}
    \item \(\frac{16}{9}\) \qty{}{m/s^2}
\end{choices}
\begin{answer}
    B\qquad 提示：\(\overline{v} = v_{\frac{t}{2}}\)
\end{answer}

\lsn{7}\icode{E-MK-SI-L2-260929-002}
运动员将足球以大小为\qty{9}{m/s}的速度踢出，足球沿草地以大小为\qty{1.5}{m/s^2}的加速度做匀减速直线运动，运动员在将足球踢出的同时以大小为\qty{6}{m/s}的速度匀速追赶足球。运动员追上足球所用的时间为\paren[B]
\begin{choices}
	\item \qty{2}{s}
	\item \qty{4}{s}
	\item \qty{6}{s}
	\item \qty{7}{s}
\end{choices}
\begin{answer}
	B\qquad 提示：作出足球和运动员的\(v - t\)图像

	\begin{tikzpicture}[scale=.8,>=stealth,font=\scriptsize]
		\draw[->,blue] (0,0)node[below left]{\(O\)}--(0,3.3)node[right]{\(v\)/(\qty{}{m\cdot s^{-1}})};
		\draw[->,blue] (0,0)--(5,0)node[below]{\(t\)/\qty{}{s}};
        \draw[dashed,cyan] (1.4,1.8)--(1.4,0) (2.8,1.8)--(2.8,0);
		\draw[purple] (0,2.7)--(4.2,0) (0,1.8)--(4.2,1.8);
		\foreach \c/\an/\n in{(0,2.7)/left/9,(0,1.8)/left/6,(4.2,0)/below/6,(2.8,0)/below/4,(1.4,0)/below/2}
			\node[blue] at \c[\an]{\(\n\)};
	\end{tikzpicture}
\end{answer}

\lsn{8}\icode{E-MK-MU-L2-260627-001}
甲、乙两辆车（均可视为质点）在同一平直道路上的两并排车道上同向行驶，初始时甲车在乙车前方\(x_0\)处，甲的位移—时间关系和乙的速度—时间关系分别为\(x_\text{甲}=t^2+10t\)和\(v_\text{乙} = t + 20\)，式中各物理量均取国际单位制单位，下列说法正确的是\paren[]
\begin{choices}
    \item 无论\(x_0\)取何值，两车都不可能相遇
    \item 若\(x_0 =\) \qty{50}{m}，两车只相遇1次
    \item 若\(x_0 =\) \qty{40}{m}，两车只相遇2次
    \item 若\(x_0 =\) \qty{10}{m}，两车只相遇1次
\end{choices}
\begin{answer}
    BC
\end{answer}

\lsn{9}\icode{E-MK-SI-L2-260929-001}
甲、乙两物体受水平力的作用，在光滑水平面上沿一条直线运动，其\(v-t\)图像如图所示，关于两物体的受力情况和运动情况，下列说法正确的是\paren[D]

\begin{tikzpicture}[>=stealth,font=\scriptsize]
	\draw[->,blue] (0,0)node[below left]{\(0\)}--(3.8,0)node[below]{\(t\)/\qty{}{s}};
	\draw[->,blue] (0,0)--(0,3)node[right]{\(v\)/(\qty{}{m\cdot s^{-1}})};
	\foreach \x/\n in {.8/2,1.6/4,2.4/6,3.2/8}
		\draw[blue] (\x,.1)--(\x,0)node[below]{\(\n\)};
	\foreach \y/\n in {.6/1,1.2/2,1.8/3,2.4/4}
		\draw[blue] (.1,\y)--(0,\y)node[left]{\(\n\)};
	\draw[purple] (0,2.4)--(3.2,0) (0,.6)--(1.6,1.2)--++($1.2*($(1.6,1.2)-(0,0.6)$)$);
	\draw[dashed,cyan] (0,1.2)--(1.6,1.2)--(1.6,0);
	\node at(.8,1.8)[above]{甲};
	\node at(2.7,1.8)[above]{乙};
\end{tikzpicture}
\begin{choices}
	\item 甲、乙两物体运动方向相反
	\item 若甲、乙两物体所受外力之比\(F_{\text{甲}}:F_{\text{乙}}=2:3\)，则质量之比\(m_{\text{甲}}:m_{\text{乙}}= 3:4\)
	\item 若\(t = 0\)时甲、乙在同一点，则\(t =\) \,\qty{4}{s}时二者相遇
	\item 若\(t =\) \,\qty{4}{s}时甲、乙相遇，则\(t =\) \,\qty{8}{s}时两物体间距离与\(t = 0\)时相等
\end{choices}
\begin{answer}
	D\quad 提示：
\end{answer}

\lsn{10}\icode{E-MK-CA-L3-260929-003}
分拣机器人在智能系统的调度下，能够自主规划路线，确保高效、准确的分拣作业。某次分拣过程示意图如图所示，机器人从\(A\)点由静止出发先沿直线\(AB\)运动到\(B\)点，此时机器人在\(B\)点的速率为零，再沿直线\(BC\)运动到\(C\)点停下，再将货物从托盘卸到分拣口。已知机器人最大运行速率 \(v_m =\) \,\qty{3}{m/s}，机器人在加速或减速运动的过程中加速度大小均为 \(a =\) \,\qty{2.5}{m/s^2}，\(A\)、\(B\)间的距离 \(x_1 =\) \,\qty{6}{m}，\(B\)、\(C\)间的距离 \(x_2 =\) \,\qty{2.5}{m}。若机器人在最短时间内到达分拣口，求：
\begin{enumerate}
	\item 机器人从\(A\)点运动到\(B\)点的过程中，由静止加速到最大运行速率 \(v_m\) 时发生的位移以及所需的时间；
	\item 机器人从\(A\)点运动到\(B\)点的时间；
	\item 机器人从\(B\)点运动到\(C\)点的平均速度大小。
\end{enumerate}

\begin{tikzpicture}[font=\scriptsize]
	\draw[dashed,gray] (0,0)coordinate(A)--(0,-3)coordinate(B)--(1.25,-3)coordinate(C);
	\foreach \x in {A,B}
		\fill (\x)node[left]{\(\x\)} circle(.05);
	\fill (C)node[below]{\(C\)} circle(.05);
	\draw (-.6,.6) rectangle(4,-3.6) (2.4,-3.1) rectangle(3.4,-2.5);
	\draw[fill=red!20] (0,0) rectangle(1,.4) (1.25,-3) rectangle(2.25,-2.6);
	\node at(.5,.2){机器人};
	\node at(1.75,-2.8){机器人};
	\node at(3.3,.2){分拣平台};
	\node at(2.9,-2.8){分拣口};
\end{tikzpicture}
\begin{answer}
	（1）\qty{1.8}{m}\quad \qty{1.2}{s}（2）\qty{3.2}{s}（3）\qty{1.25}{m/s}

	机器人从\(A\)到\(B\) 和从\(B\)到\(C\)的运动图像如下图所示

	\begin{tikzpicture}[>=stealth,font=\scriptsize]
		\draw[->,blue] (0,0)node[below left]{0}--(4,0)node[below]{\(t\)/\qty{}{s}};
		\draw[->,blue] (0,0)--(0,3)node[right]{\(v\)/(\qty{}{m\cdot s^{-1}})};
		\draw[dashed,cyan] (0,2)node[blue,left]{3}--(1.2,2)--(1.2,0)node[blue,below]{1.2} (2,2)--(2,0)node[blue,below]{2} (0,{5/3})node[blue,left]{2.5}--(1,{5/3})--(1,0);
		\draw[purple] (0,0)--(1.2,2)--(2,2)--(3.2,0)node[blue,below]{3.2};
		\draw[teal,thin] (0,0)--(1,{5/3})--(2,0);
	\end{tikzpicture}
\end{answer}
\vspace*{3cm}

\lsn{11}\icode{N13-MK-CA-L3-260928-002}
\nceeID[13]{2013课标\Rmnum{1}}{24}
水平桌面上有两个玩具车\(A\)和\(B\)，两者用一轻质细橡皮筋相连，在橡皮筋上有一红色标记\(R\)。在初始时橡皮筋处于拉直状态，\(A\)、\(B\)和\(R\)分别位于直角坐标系中的(0,\(2l\))、(0,\( - l\))和(0,0)点。已知\(A\)从静止开始沿\(y\)轴正向做加速度大小为\(a\)的匀加速运动；\(B\)平行于\(x\)轴朝\(x\)轴正向匀速运动。在两车此后运动的过程中，标记\(R\)在某时刻通过点(\(l\),\(l\))。假定橡皮筋的伸长是均匀的，求\(B\)运动速度的大小。
\begin{answer}
    \(v = \frac{1}{4}\sqrt{6al}\)
\end{answer}